% -----------------------------------------------------------------------------------------------------
\section{规范不变性、自由场值与支撑}\label{sec:pdf-properties}

\begin{proposition}[规范不变性、推动不变性与自由场值 \tagR]\label{prop:pdf-invariance}
（$R_{\mathrm{def}}$：由算符定义直接推出。）
(a) $f_{q/H}$、$f_{\bar q/H}$、$f_{g/H}$ 在规范变换 $\psi\to U\psi$，$A_\mu\to UA_\mu U^\dagger-\frac{\ii}{g}(\partial_\mu U)U^\dagger$ 下不变。
(b) 它们只依赖于 $\xi$（与 $\mu$），在沿 $z$ 轴的推动 $v^\pm\to\ee^{\pm\eta}v^\pm$ 下不变。
(c) 对动量 $p=(p^+,0^-,\bm 0_{\mathrm T})$ 的无质量自由靶，在 $g^0$ 阶
\begin{equation}\label{eq:pdf-free}
f^{(0)}_{q/q}=\delta(1-\xi),\qquad f^{(0)}_{q/\bar q}=-\delta(1+\xi)\ \ \big(\text{故 }f^{(0)}_{\bar q/\bar q}=\delta(1-\xi)\big),\qquad
f^{(0)}_{g/g}=\delta(1-\xi)-\delta(1+\xi).
\end{equation}
% VF-PDF-1, VF-PDF-2
(d) 胶子的归一由能动张量固定：$T^{++}_g=\sum_jG^{+j}_aG^{+j}_a$ 满足 $\langle k|T^{++}_g|k\rangle=2(k^+)^2$，与 $\langle P|T^{\mu\nu}|P\rangle=2P^\mu P^\nu$ 一致。
非微扰地位：强子态需\cref{asm:intro-wightman}、\cref{asm:intro-composite}（MA0）。
\uses{def:pdf-quark,def:pdf-antiquark,def:pdf-gluon,prop:prelim-wilson-covariance,lem:prelim-boost,asm:intro-wightman,asm:intro-composite}
\end{proposition}

\begin{proof}
(a) 由\cref{prop:prelim-wilson-covariance}，$W[w,0]\to U(w)W[w,0]U^\dagger(0)$，故 $\bar\psi(w)\gamma^+W[w,0]\psi(0)$ 不变。胶子：记 $G=G_at^a$，$G_a(w)W^{(A)}_{ab}[w,0]G_b(0)=2\Tr[G(w)W[w,0]G(0)W[0,w]]$，而 $G\to UGU^\dagger$，迹中的 $U$ 两两相消。
(b) 由\cref{lem:prelim-boost}，好分量 $P_+\psi$（$P_+=\frac12\gamma^-\gamma^+$）按 $\ee^{\eta/2}$ 标度；$\gamma^0\gamma^+=\sqrt2P_+$ 给出 $\bar\psi\gamma^+\psi=\sqrt2\,\psi_+^\dagger\psi_+$，故双线性式按 $\ee^{\eta}$ 标度，而 $w^-\to\ee^{-\eta}w^-$。于是 $M_P(w^-)\equiv\langle P|\bar\psi(w)\gamma^+W[w,0]\psi(0)|P\rangle_{\mathrm c}$ 满足 $M_{\Lambda P}(w^-)=\ee^{\eta}M_P(\ee^{\eta}w^-)$（协变归一在推动下不变），代入\cref{eq:pdf-fq}并换元 $w'=\ee^\eta w^-$ 即得 $f$ 不变。胶子同理（$G^{+j}G^{+j}$ 按 $\ee^{2\eta}$ 标度，$1/(\xi P^+)$ 补偿一个 $\ee^{\eta}$）。
(c) 在 $g^0$ 阶 $W=1$，$\psi=\int[b\,u\,\ee^{-\ii p\cdot x}+d^\dagger v\,\ee^{\ii p\cdot x}]$。夸克靶的连通矩阵元为 $\bar u\gamma^+u\,\ee^{\ii p^+w^-}=2p^+\ee^{\ii p^+w^-}$（对每个螺旋度 $\bar u\gamma^+u=2p^+$，显式验证见下例），故 $f^{(0)}_{q/q}=\frac{2p^+}{4\pi}\,2\pi\delta((1-\xi)p^+)=\delta(1-\xi)$。反夸克靶：$\bar\psi(w)$ 中的 $d\,\bar v\,\ee^{-\ii p\cdot w}$ 与 $\psi(0)$ 中的 $d^\dagger v$ 缩并，化为正规序时 $dd^\dagger\to-d^\dagger d$，矩阵元为 $-\bar v\gamma^+v\,\ee^{-\ii p^+w^-}=-2p^+\ee^{-\ii p^+w^-}$，故 $f^{(0)}_{q/\bar q}=-\delta(1+\xi)$。胶子靶（$\bm k_{\mathrm T}=0$）：$\langle0|G^{+j}_a(x)|k,\lambda,b\rangle=-\ii k^+\epsilon^j_\lambda\delta_{ab}\ee^{-\ii k\cdot x}$；连通矩阵元含“$0$ 处湮灭、$w$ 处产生”及其相反两项，对极化平均后 $\sum_j|\epsilon^j|^2\to1$，得 $(k^+)^2(\ee^{\ii k^+w^-}+\ee^{-\ii k^+w^-})$，于是 $f^{(0)}_{g/g}=\frac{(k^+)^2}{2\pi\xi k^+}\frac{2\pi}{k^+}[\delta(\xi-1)+\delta(\xi+1)]=\delta(1-\xi)-\delta(1+\xi)$（$\xi=-1$ 处 $1/\xi=-1$）。
(d) 同一胶子态中 $\sum_jG^{+j}G^{+j}$ 的矩阵元是上式在 $w^-=0$ 处的值 $2(k^+)^2$；另一方面 $T^{++}_g=-G^{+\alpha}_aG^{+}{}_{\alpha a}$ 是规范不变能动张量 $-G^{\mu\alpha}_aG^\nu{}_{\alpha a}+\frac14g^{\mu\nu}G^2$ 的 $++$ 分量（$G^{++}=0$，$G^+{}_-=0$，只剩横向求和），$\langle k|T^{\mu\nu}|k\rangle=2k^\mu k^\nu$ 给出同一值。
\end{proof}

\begin{example}[自由场值]\label{ex:pdf-free-field}
Dirac 表示中，$p$ 沿 $+z$、$E=p^+/\sqrt2$，螺旋度 $h$ 的旋量 $u=\sqrt E(\chi_h,2h\chi_h)^{\tp}$；由 $\gamma^0\gamma^+=(1+\gamma^0\gamma^3)/\sqrt2$，$\bar u\gamma^+u=\frac{E}{\sqrt2}(2+2(2h)^2)=2p^+$，与 $h$ 无关。割图表示\cref{eq:pdf-cut}中夸克靶只有 $X=$ 真空一项，给出 $\frac12\delta(\xi p^+-p^+)\cdot2p^+=\delta(1-\xi)$；反夸克靶的负号来自费米统计而非旋量（用\cref{def:pdf-antiquark}中电荷共轭场的定义直接得 $f^{(0)}_{\bar q/\bar q}=\delta(1-\xi)$）。各自由分布的矩 $1$、$(-1)^n$、$1+(-1)^n$ 与定域算符 $\bar\psi\gamma^+(\ii\partial^+)^{n-1}\psi$、$\sum_jG^{+j}(\ii\partial^+)^{n-2}G^{+j}$ 在同样态中的值除以 $2(p^+)^n$ 逐一相等（\cref{thm:pdf-moments-local}的自由场特例）；自由胶子的 $\int_0^1x\,g=1$ 即动量求和规则（\cref{prop:pdf-sum-rules}）。
\uses{prop:pdf-invariance}
\end{example}

\begin{theorem}[支撑（微扰论）\tagR]\label{thm:pdf-support-pt}
（$R_{\mathrm{PT}}$：在重整化微扰论中对红外正规化的部分子态严格。）
设靶为红外正规化的在壳部分子 $j$。在微扰论每一阶、对每个割图，裸分布 $f_{q/j}(\xi)$ 与 $f_{g/j}(\xi)$ 在 $|\xi|>1$ 时为零，即它们是支撑在 $[-1,1]$ 上的分布（等价地，$x>0$ 处的各密度支撑在 $[0,1]$ 上）。$\MSbar$ 重整化保持这一性质，因为逐阶由极点部分构造的抵消项 $Z_{ij}(z)$ 支撑在 $0\le z\le1$ 上（$Z(\xi/\zeta)$ 在 $\xi>\zeta$ 时为零）。（本定理只涉及部分子态。）
\uses{def:pdf-quark,def:pdf-gluon,def:ir-landau-setup}
\end{theorem}

\begin{proofsketch}
(1) 微扰论中\cref{eq:pdf-fq}的矩阵元是普通（非编时）乘积；在 $[\bar\psi W](w)$ 与 $[W\psi](0)$ 之间插入末态完备集得割图求和\cref{eq:pdf-cut}。每个割图中穿过割线的都是在壳正能量粒子，各自 $p^+=(E+p^3)/\sqrt2\ge0$，故 $p_X^+\ge0$，$\delta$ 函数给出 $\xi p^+=p^+-p_X^+\le p^+$。Wilson 线只是源，不携带在壳末态。(2) 由\cref{eq:pdf-cconj}，$f_q(\xi)=-f_{\bar q}(-\xi)$，对电荷共轭场的分布重复 (1) 得 $\xi\ge-1$；反对易子的 c 数与 $\xi$ 无关，不进入连通矩阵元。(3) 胶子以伴随 Wilson 线同样处理。(4) 对部分子靶，$Z^{(k)}=-\mathrm{PP}\big[f^{(k)}_{\mathrm{bare}}+\sum_{l<k}Z^{(l)}\otimes f^{(k-l)}_{\mathrm{bare}}\big]$（PP 为 $\varepsilon$ 的极点部分）；支撑在 $[0,1]$ 上的分布的 Mellin 卷积仍支撑在 $[0,1]$ 上，取 Laurent 系数不扩大支撑，归纳即得。这是谱条件论证在图上的版本；用光锥序微扰论的表述见\cite[第~6.9.3~节，第~191~页]{Collins2011}（该处指出论证是非微扰的），对多部分子关联函数的推广见\cite{Jaffe1983}（据该文摘要，每个量子的动量分数都在 $0$ 与 $1$ 之间），时间排序与支撑的关系见\cite{Diehl:1998sm}（据该文摘要）。完整论证见上述文献。
\end{proofsketch}

\begin{proposition}[支撑（强子态）\tagA]\label{prop:pdf-support-hadron}
对强子态，$f_{q/H}(\xi)=0$ 与 $f_{g/H}(\xi)=0$ 对 $|\xi|>1$ 成立：(a) 在 $\bar\psi(w)W$ 与 $W\psi(0)$ 之间插入完备态，$w^-$ 积分把中间态的加号动量固定为 $p_X^+=(1-\xi)P^+$，谱条件给出 $p_X^+\ge0$，故 $\xi\le1$；(b) 反夸克部分由光类间隔上的反对易子给出（扣除与态无关的发散常数），由此 $\xi\ge-1$。
非微扰地位：\cref{asm:intro-wightman}、\cref{asm:intro-composite}（MA0）。
\uses{def:pdf-quark,def:pdf-antiquark,thm:pdf-support-pt,asm:intro-wightman,asm:intro-composite}
\end{proposition}

\begin{proofsketch}
(1) 由\cref{asm:intro-wightman}(W2)，联合谱在闭前向锥内，故任何中间态 $p_X^0\ge|\bm p_X|$，$p_X^+\ge0$。(2) 插入 $\sum_X|X\rangle\langle X|$，与\cref{thm:pdf-support-pt}第 (1) 步相同地得 $\xi\le1$。(3) 写 $\bar\psi(w)\Gamma\psi(0)=-\psi(0)^{\tp}\Gamma^{\tp}\bar\psi^{\tp}(w)+\{\cdot,\cdot\}$；反对易子在光类间隔上是与态无关的发散常数，从连通矩阵元中扣除；对剩下的项插入完备态得 $\xi\ge-1$。(4) 胶子同理。论证是非微扰的\cite[第~6.9.3~节，第~191~页]{Collins2011}，微扰论中的对应是\cref{thm:pdf-support-pt}。
\end{proofsketch}

\paragraph{地位与缺口}
谱论证把\emph{重整化的}光线算符当作可以在两半之间插入完备态的算符乘积；而重整化作用在整个双局域算符上，并不是两个“半算符”各自重整化之积。“对态求和”与“重整化”的交换是未经证明的一步，它使本命题停留在 \tagA；所依赖的假设是\cref{asm:intro-wightman}(W2) 与\cref{asm:intro-composite}(C3)（光线算符在强子态之间作为缓增分布存在），这一步隐含在\cite[第~6.9.3~节]{Collins2011}的谱论证中。此缺口与幂次计数无关，不涉及 E-S 至 E-W 型的误差估计。$\xi=0$ 一点需要小心：扭度 3 的分布可以含 $\delta(x)$ 项，目前不知道扭度 2 的非极化 PDF 有这类项。若能非微扰地构造重整化光线算符并证明其谱表示，结论即可在\cref{asm:intro-wightman}、\cref{asm:intro-composite}之下升级为 $[R,R_{\mathrm{cond}}]$。第二种论证见\cref{rem:pdf-support-jaffe}。

\begin{remark}[支撑的第二种论证：$k^-$ 极点]\label{rem:pdf-support-jaffe}
写 $f(\xi)=\int\frac{\dd^4k}{(2\pi)^4}\delta(k^+-\xi p^+)\TrD[\frac12\gamma^+\mathcal A(k,p)]$，$\mathcal A$ 是部分子—靶前向振幅。携带加号动量 $\ell^+$ 的传播子在 $\ell^-=(\bm\ell_{\mathrm T}^2+m^2-\ii0)/(2\ell^+)$ 处有极点，$\ell^+>0$ 时在实轴下方，$\ell^+<0$ 时在上方。当 $\xi>1$ 时，靶取走一个部分子后的中间态总加号动量 $(1-\xi)p^+$ 为负（$\xi<-1$ 时交换夸克与反夸克的角色），完成其余回路的减号分量积分后，$\mathcal A$ 作为 $k^-$ 的函数的奇点全在实轴同一侧；被积函数在 $|k^-|\to\infty$ 时足够快地衰减，向另一侧闭合围道即得零\cite{Jaffe1983}（据该文摘要，论证要求中间强子态在壳且没有相对于靶或中间态半不连通的图）。微扰论中它与\cref{thm:pdf-support-pt}等价，两者都表述“中间态加号动量非负”。
\uses{thm:pdf-support-pt}
\end{remark}

% -----------------------------------------------------------------------------------------------------
\section{数密度诠释及其局限（启发式）；光锥规范的剩余规范自由}\label{sec:pdf-interpretation}

算符定义的直观读法来自光锥规范 $A^+=0$（此时 $W[w,0]=1$）；它只是启发式的，并且正是“$\MSbar$ PDF 是概率、故为正”这一误解的来源。

\begin{remark}[数密度诠释（启发式）]\label{rem:pdf-number-density}
在 $A^+=0$ 规范中按光前产生、湮灭算符（$\{b,b^\dagger\}=(2\pi)^32k^+\delta^{(3)}$）展开好分量，用 $\bar\psi\gamma^+\psi=\sqrt2\psi_+^\dagger\psi_+$、$u_+^\dagger u_+=\sqrt2k^+$，对 $0<\xi<1$ 得
\begin{equation}\label{eq:pdf-number}
f_{q/H}(\xi)=\sum_{s,c}\int\frac{\dd^{2}\bm k_{\mathrm T}}{2\xi\,(2\pi)^3}\,
\frac{\big\|b_{s,c}(\xi P^+,\bm k_{\mathrm T})|P\rangle\big\|^2}{\langle P|P\rangle}\ \ge0
\end{equation}
（对自由夸克，比值为 $(2\pi)^32p^+\delta^{(3)}(k-p)$，给出 $\delta(1-\xi)$），即“份额为 $\xi$ 的夸克数密度”\cite[第~6.7~节、第~6.9.1~节]{Collins2011}。然而：(i) 相互作用理论中 $\bm k_{\mathrm T}$ 积分在紫外对数发散，正是 $\MSbar$ 所要去掉的发散，重整化后正性不复存在——Collins 写道：「“this particular result will not hold exactly in renormalizable theories, because of the need for renormalization”」\cite[第~6.9.1~节]{Collins2011}，即因为需要重整化，该结果在可重整理论中不严格成立；(ii) $A^+=0$ 规范的伪 $1/k^+$ 奇异性需要处方（\cref{rem:pdf-gauge-residual}）；(iii) 光前量子化有 $k^+=0$ 的零模问题；(iv) 光锥无穷远处的横向规范连接在 $A^+=0$ 规范中并不消失\cite{JiYuan2002,BelitskyJiYuan2003}，对共线 PDF 它们因两端横向位置相同而相消，对横动量依赖分布则不然。所以\cref{eq:pdf-number}是关于\emph{裸的、紫外截断的}量的启发式陈述，不是关于 $\MSbar$ PDF 的陈述；后者的支撑与求和规则另有证明（\cref{thm:pdf-support-pt}、\cref{prop:pdf-sum-rules}），正定性是\cref{conj:pdf-positivity}。
\end{remark}

\begin{remark}[光锥规范的剩余规范自由与 $1/k^+$ 处方]\label{rem:pdf-gauge-residual}
$A^+=0$ 并未完全固定规范：保持它的变换与 $x^-$ 无关，作用在 $x^-=\pm\infty$ 处；相应地，传播子在 $k^+=0$ 的奇异性需要处方（主值、$k^+\pm\ii0$、Mandelstam--Leibbrandt 等），每种处方等价于对 $\bm A_{\mathrm T}$ 在 $x^-=\pm\infty$ 的一种边界条件。共线 PDF 的 Wilson 线是横向位置相同的\emph{有限}类光线段，$W[w,0]=W[w,\pm\infty n_B]W[\pm\infty n_B,0]$，无穷远横向连接平凡，故结果与处方无关；横动量依赖分布的订书钉形 Wilson 线两端横向位置不同，结果依赖于边界条件（\cref{ch:tmd}）。常见论证“用规范变换把 $W$ 变成 $1$”中的 $U(x)=W[\infty n_B,x]$ 是依赖于场的算符而非整体规范变换；结论仍成立，只因它用到的仅是合成律。
\end{remark}
